\chapter*{Abstract}\label{abstract}
\phantomsection
\pdfbookmark[1]{Abstract}{abstract}%
\prevspace
This thesis develops a two-stage electrocardiogram (ECG) system that uses one
lead. The first supervised stage decides whether a heartbeat is normal or
abnormal. If the beat is abnormal, the second stage classifies it as premature
ventricular contraction (PVC), premature atrial contraction (PAC), left bundle
branch block (LBBB), right bundle branch block (RBBB), or atrial fibrillation
(AFib).

All models receive one ECG channel at a time. A dilated U-Net autoencoder
learned normal ECG structure from the complete first stored ECG channel of all
18 MIT--BIH Normal Sinus Rhythm Database records. The WFDB header names this
channel ECG1, but it does not state the exact anatomical placement. Fifteen
subjects were used for fitting and three different subjects were used for
validation. Four reconstruction-error measures were then combined with MLII
beat shape, frequency power, and causal RR-interval context. Two separate Extra
Trees ensembles performed abnormality detection and abnormal-subtype
classification.

The supervised study used every MIT--BIH Arrhythmia Database record that
contains MLII: 46 records representing 45 subject clusters. In each record, the
first 80\% of time was used for development and the final 20\% was kept for
testing. A 16-beat gap was left before the test boundary to reduce contact
between neighbouring windows. Therefore, the test contains later unseen signal
from known patients. It is not an unseen-patient test.

On 20,988 temporal-test beats, abnormality accuracy was 98.47\%, balanced accuracy was
98.22\%, sensitivity was 96.59\%, and specificity was 99.85\%. Conditional
subtype accuracy on 7,964 supported abnormal beats was 95.20\%, with macro F1
of 0.9325. On 20,062 supported six-class beats, whole-system accuracy was
96.86\% and macro F1 was 0.9303. Equal-subject whole-system accuracy was
96.90\%. A 10,000-replicate subject-cluster bootstrap gave a 95\% confidence
interval of 94.02--99.05\%. Its lower bound exceeds 90\%, but its 5.03-point
width does not meet the planned three-point precision target. Forty-two of 45
subject clusters exceeded 90\% ($p=4.33\times10^{-10}$, exact one-sided sign
test).

The aggregate target was reached for this known-patient temporal holdout, but
not for every subtype. PAC end-to-end recall was 70.05\%. The system therefore
shows strong performance on later ECG after patient-specific training. However,
it does not prove more than 90\% accuracy for completely new patients or for
every arrhythmia type. It remains a research prototype, not a clinical device.
